% English translation of chapter2.tex lines 1564--1754.
% Source: Methods of Algebra, Volume 2, commit
% 9a5803ff77dd3257484cb177f851a73770a59dd3 (CC BY 4.0).
% stable-unit-id: o014.aljabr2.chapter2.serre-subcategories-and-k0-groups
% This segment stops exactly before the source section at line 1756.

% segment-id: o014.li2.chapter2.serre-subcategories-and-k0-groups.h001
\section{Serre Subcategories and \texorpdfstring{$\mathrm{K}_0$}{K0} Groups}\label{sec:Serre-subcat}
\hypertarget{o014.aljabr2.chapter2.serre-subcategories-and-k0-groups}{ }

% segment-id: o014.li2.chapter2.serre-subcategories-and-k0-groups.p001
A full subcategory of an abelian category automatically inherits the structure
of an $\cate{Ab}$-category, so one may ask whether the full subcategory is an
additive category or an abelian category.

% segment-id: o014.li2.chapter2.serre-subcategories-and-k0-groups.d001
\begin{definition}\index{abelian category!abelian subcategory}
If $\mathcal{B}$ is a full subcategory of an abelian category $\mathcal{A}$,
$\mathcal{B}$ is itself abelian, and the inclusion functor
$\iota:\mathcal{B}\to\mathcal{A}$ is exact
(Definition~\ref{def:exact-functor}), then $\mathcal{B}$ is called an
\emph{abelian subcategory} of $\mathcal{A}$.
\end{definition}

% segment-id: o014.li2.chapter2.serre-subcategories-and-k0-groups.pr001
\begin{proposition}\label{prop:Abel-subcat}
A full subcategory $\mathcal{B}$ of an abelian category $\mathcal{A}$ is an
abelian subcategory if and only if the following conditions hold.
% segment-id: o014.li2.chapter2.serre-subcategories-and-k0-groups.l001
\begin{itemize}
% segment-id: o014.li2.chapter2.serre-subcategories-and-k0-groups.i001
	\item $0\in\Obj(\mathcal{B})$;
% segment-id: o014.li2.chapter2.serre-subcategories-and-k0-groups.i002
	\item if $X,Y\in\Obj(\mathcal{B})$, then $X\oplus Y$ may also be chosen as
	an object of $\mathcal{B}$;
% segment-id: o014.li2.chapter2.serre-subcategories-and-k0-groups.i003
	\item for every morphism $f:X\to Y$, if
	$X,Y\in\Obj(\mathcal{B})$, then $\Ker(f)$ and $\Coker(f)$ may also be
	chosen as objects of $\mathcal{B}$.
\end{itemize}
\end{proposition}

% segment-id: o014.li2.chapter2.serre-subcategories-and-k0-groups.pf001
\begin{proof}
Notice that these conditions are self-dual. The ``only if'' direction is
clear. Now suppose that the conditions hold. Then $\mathcal{B}$ is an additive
category and $\iota:\mathcal{B}\to\mathcal{A}$ is an additive functor. We next
show that $\iota$ preserves all finite $\varprojlim$ and $\varinjlim$. First,
$\iota$ preserves $0$ and finite direct sums. Moreover, equalizers can be
expressed as $\Ker$ (Remark~\ref{rem:equalizer-Ker}), while the stated
condition says that $\mathcal{B}$ is closed under taking $\Ker$. Thus $\iota$
preserves all finite $\varprojlim$; the case of $\varinjlim$ is dual.

Recalling the definitions of $\Image$ and $\Coim$ in
Definition~\ref{def:Im-Coim} and using the preceding step, we see that
$\mathcal{B}$ is closed under both constructions. For every morphism $f$ in
$\mathcal{B}$, the canonical morphism
$\Coim(f)\rightiso\Image(f)$ in diagram \eqref{eqn:strict-morphism} is an
isomorphism in $\mathcal{A}$, hence also in $\mathcal{B}$. Therefore
$\mathcal{B}$ is an abelian subcategory.
\end{proof}

% segment-id: o014.li2.chapter2.serre-subcategories-and-k0-groups.d002
\begin{definition}[J.-P.\ Serre]\label{def:Serre-subcat}
\index{Serre subcategory}
\index{weak Serre subcategory}
A full subcategory $\mathcal{T}$ of an abelian category $\mathcal{A}$ is
called a \emph{Serre subcategory} of $\mathcal{A}$ if it satisfies the
following conditions.
% segment-id: o014.li2.chapter2.serre-subcategories-and-k0-groups.l002
\begin{itemize}
% segment-id: o014.li2.chapter2.serre-subcategories-and-k0-groups.i004
	\item $0\in\Obj(\mathcal{T})$;
% segment-id: o014.li2.chapter2.serre-subcategories-and-k0-groups.i005
	\item for every short exact sequence
	$0\to X'\to X\to X''\to0$ in $\mathcal{A}$, one has
	$X\in\Obj(\mathcal{T})$ if and only if
	$X',X''\in\Obj(\mathcal{T})$.
\end{itemize}
If the last condition is weakened as follows: for every exact sequence in
$\mathcal{A}$
% segment-id: o014.li2.chapter2.serre-subcategories-and-k0-groups.q001
\[W\to X'\to X\to X''\to Y,\]
if $W,X',X'',Y\in\Obj(\mathcal{T})$, then $X\in\Obj(\mathcal{T})$, then
$\mathcal{T}$ is called a \emph{weak Serre subcategory} of $\mathcal{A}$
\footnote{This somewhat abrupt definition is designed specifically for
derived categories; see \S\sourcecrossref{sec:derived-cat}{4.4}.}.
\end{definition}

% segment-id: o014.li2.chapter2.serre-subcategories-and-k0-groups.p002
For example, if $R$ is a commutative ring, all Noetherian (or Artinian)
modules form a Serre subcategory of $R\dcate{Mod}$. This is a standard fact in
module theory \cite[Lemma 6.10.2]{Li1}.

% segment-id: o014.li2.chapter2.serre-subcategories-and-k0-groups.p003
If $\mathcal{T}$ is a Serre subcategory (or weak Serre subcategory) of
$\mathcal{A}$, then $\mathcal{T}^{\opp}$ is a Serre subcategory (or weak
Serre subcategory) of $\mathcal{A}^{\opp}$.

% segment-id: o014.li2.chapter2.serre-subcategories-and-k0-groups.p004
A weak Serre subcategory $\mathcal{T}$ has the following saturation property:
if $X\in\Obj(\mathcal{A})$ is isomorphic to an object of $\mathcal{T}$, then
$X\in\Obj(\mathcal{T})$. This follows by applying the last condition to the
exact sequence $0\to0\to X\rightiso Y\to0$.

% segment-id: o014.li2.chapter2.serre-subcategories-and-k0-groups.co001
\begin{corollary}\label{prop:Serre-subcat-abelian}
If $\mathcal{T}$ is a weak Serre subcategory of an abelian category
$\mathcal{A}$, then $\mathcal{T}$ is an abelian subcategory of $\mathcal{A}$.
\end{corollary}

% segment-id: o014.li2.chapter2.serre-subcategories-and-k0-groups.pf002
\begin{proof}
It suffices to verify the conditions of Proposition~\ref{prop:Abel-subcat}.
First, $0\in\Obj(\mathcal{T})$. Next, insert the following exact sequences
into the definition of a weak Serre subcategory:
% segment-id: o014.li2.chapter2.serre-subcategories-and-k0-groups.g001
\begin{gather*}
	0\to X\to X\oplus Y\to Y\to0
	\quad\text{(Proposition~\ref{prop:biproduct-ses})},\\
	0\to0\to\Ker(f)\to X\xrightarrow{f}Y,\quad
	X\xrightarrow{f}Y\to\Coker(f)\to0\to0,
\end{gather*}
which shows that $\mathcal{T}$ is closed under direct sums, $\Ker$, and
$\Coker$.
\end{proof}

% segment-id: o014.li2.chapter2.serre-subcategories-and-k0-groups.ex001
\begin{example}\index[sym1]{kerF@$\Ker(F)$}
If $F:\mathcal{A}\to\mathcal{B}$ is an exact functor between abelian
categories, then the objects satisfying $FX=0$ form a Serre subcategory of
$\mathcal{A}$, denoted by $\Ker(F)$.
\end{example}

% segment-id: o014.li2.chapter2.serre-subcategories-and-k0-groups.th001
\begin{theorem}[Serre quotient]\label{prop:Serre-quotient}
\index{Serre quotient}
Let $\mathcal{T}$ be a Serre subcategory of an abelian category
$\mathcal{A}$. Then there is an abelian category
$\mathcal{A}/\mathcal{T}$ (which may be a large category), together with an
exact, essentially surjective functor
$Q:\mathcal{A}\to\mathcal{A}/\mathcal{T}$, such that
$\Ker(Q)=\mathcal{T}$ and with the following universal property: for every
abelian category $\mathcal{B}$ and every exact functor
$F:\mathcal{A}\to\mathcal{B}$ with $\Ker(F)\supset\mathcal{T}$, there is a
unique exact functor $G:\mathcal{A}/\mathcal{T}\to\mathcal{B}$ such that
$F=GQ$.

In this universal property, $G$ is faithful if and only if
$\Ker(F)=\mathcal{T}$.
\end{theorem}

% segment-id: o014.li2.chapter2.serre-subcategories-and-k0-groups.pf003
\begin{proof}
Define
$S:=\left\{f\in\Mor(\mathcal{A}):\Ker(f),\Coker(f)\in\Obj(\mathcal{T})\right\}$.
We shall show that $S$ is a multiplicative system in the sense of
Definition~\ref{def:multiplicative-family}.

Clearly $S$ contains all identity morphisms, so (S1) holds. Let $f:X\to Y$
and $g:Y\to Z$ belong to $S$. The mutually dual exact sequences
% segment-id: o014.li2.chapter2.serre-subcategories-and-k0-groups.g002
\begin{gather*}
	0\to\Ker(f)\to\Ker(gf)\xrightarrow{f}\Ker(g),\\
	\Coker(f)\xrightarrow{g}\Coker(gf)\to\Coker(g)\to0,
\end{gather*}
immediately give $gf\in S$, so (S2) holds. Next consider morphisms
$X\xrightarrow{s\in S}Z\xleftarrow{f}Y$. Define
$W:=X\dtimes{Z}Y$, with morphism $s':W\to Y$. Observe that
$\Ker(s')\simeq\Ker(s)$ (Proposition~\ref{prop:pull-back-ker}), and $f$
induces $\Coker(s')\hookrightarrow\Coker(s)$
(Corollary~\ref{prop:Abel-cat-Coker-pull}); hence $s'\in S$. Thus (S3) holds.
Finally, consider morphisms
% segment-id: o014.li2.chapter2.serre-subcategories-and-k0-groups.g003
$\begin{tikzcd}
	X \arrow[r, "f", shift left] \arrow[r, "g"', shift right] &
	Y \arrow[r, "s \in S"] & W
\end{tikzcd}$,
with $sf=sg$. Since $\Image(f-g)\hookrightarrow\Ker(s)$, we have
$\Image(f-g)\in\Obj(\mathcal{T})$; therefore
$Z:=\Ker(f-g)\hookrightarrow X$ is a morphism in $S$, proving (S4).

Thus $S$ is a left multiplicative system. All the conditions are self-dual,
so $S$ is also a right multiplicative system. Now take the localization
functor
$Q:\mathcal{A}\to\mathcal{A}/\mathcal{T}:=\mathcal{A}[S^{-1}]$.
Proposition~\ref{prop:localization-extra-property} says that $Q$ is
essentially surjective (indeed,
$\Obj(\mathcal{A})=\Obj(\mathcal{A}/\mathcal{T})$), while
Proposition~\ref{prop:Abel-cat-localization} says that $Q$ is an exact functor
between abelian categories. Notice that
$\identity_{QX}=Q(\identity_X)$ equals $0$ if and only if there is a
$U\in\Obj(\mathcal{A})$ such that $U\xrightarrow{0}X$ lies in $S$ (this can
be checked directly using Corollary~\ref{prop:fraction-zero}); this implies
$X=\Coker[U\xrightarrow{0}X]\in\Obj(\mathcal{T})$. Conversely, if
$X\in\Obj(\mathcal{T})$, we may take $U=X$. Hence
$\Ker(Q)=\mathcal{T}$.

We next verify the universal property. If
$\Ker(F)\supset\mathcal{T}$, then the exact sequence
$\Ker(s)\to X\xrightarrow{s\in S}Y\to\Coker(s)$ shows that $F$ sends every
element of $S$ to an isomorphism, and the universal property of localization
determines the required $G$.

It remains to verify that $G$ is exact.
Theorem~\ref{prop:localization-additivity} first ensures that $G$ is additive.
Since every morphism in $\mathcal{A}[S^{-1}]$ can, after composition with an
isomorphism arising from $S$, be represented by a morphism in $\mathcal{A}$,
checking that $G$ preserves kernels reduces to the exactness of $F$ and $Q$.
The same applies to cokernels. Since $Q$ is essentially surjective and
$\Ker(Q)=\mathcal{T}$, the characterization of faithfulness reduces to
statement (ii) of Proposition~\ref{prop:faithful-exact}.
\end{proof}
% segment-id: o014.li2.chapter2.serre-subcategories-and-k0-groups.p005
Notice that $\mathcal{A}/\mathcal{T}$ is constructed as a localization and may
therefore be a ``large category.'' The exercises give another description of
$\mathcal{A}/\mathcal{T}$ that controls the size of the Serre quotient under
the assumption that $\mathcal{A}$ is a well-powered category
(Definition~\ref{def:well-powered}).
\index{size issue}

% segment-id: o014.li2.chapter2.serre-subcategories-and-k0-groups.ex002
\begin{example}
Let $S$ be a multiplicative subset of a commutative ring $R$. Define a full
subcategory $\mathcal{T}$ of $R\dcate{Mod}$ by declaring that
$M\in\Obj(\mathcal{T})$ if and only if, for every $m\in M$, there is an
$s\in S$ such that $sm=0$. It is easy to verify that $\mathcal{T}$ is a Serre
subcategory. We now show that $R\dcate{Mod}/\mathcal{T}$ is equivalent to
$R[S^{-1}]\dcate{Mod}$.

The assignment $M\mapsto M[S^{-1}]:=M\dotimes{R}R[S^{-1}]$ defines an exact
functor $F:R\dcate{Mod}\to R[S^{-1}]\dcate{Mod}$ that sends $\mathcal{T}$ to
zero, so the universal property gives an exact functor
$G:R\dcate{Mod}/\mathcal{T}\to R[S^{-1}]\dcate{Mod}$. Clearly
$\Ker(F)=\mathcal{T}$, so $G$ is faithful. Moreover, $G$ is essentially
surjective: regard any $R[S^{-1}]$-module $N$ as an $R$-module; one checks
directly that there is an isomorphism of $R[S^{-1}]$-modules
$N\dotimes{R}R[S^{-1}]\simeq N$.

It remains to prove that $G$ is full and faithful. Given an
$R[S^{-1}]$-module homomorphism
$\varphi:M_1[S^{-1}]\to M_2[S^{-1}]$, form the fiber product
\footnote{Editorial correction (O014-C029): the source defines $M_0$ as the
set of elements $m\in M_1$ whose images ``come from'' $M_2$, and then uses a
restriction $\varphi|_{M_0}:M_0\to M_2$. A preimage in $M_2$ is not canonical
when $M_2\to M_2[S^{-1}]$ has an $S$-torsion kernel, so that morphism is not
well-defined. We use instead the canonical fiber product and its two
projections.}
$M_0:=M_1\dtimes{M_2[S^{-1}]}M_2$, where the map
$M_1\to M_2[S^{-1}]$ is the composite
$M_1\to M_1[S^{-1}]\xrightarrow{\varphi}M_2[S^{-1}]$.
Write its projections as $p_1:M_0\to M_1$ and $p_2:M_0\to M_2$. The kernel
and cokernel of $p_1$ are $S$-torsion modules, so $p_1$ lies in the
multiplicative system defining the Serre quotient. Hence the roof
% segment-id: o014.li2.chapter2.serre-subcategories-and-k0-groups.g004
$\begin{tikzcd}
	M_1 & M_0 \arrow[l, "p_1"'] \arrow[r, "p_2"] & M_2
\end{tikzcd}$
defines a morphism in $R\dcate{Mod}/\mathcal{T}$ that is sent to $\varphi$.
This proves the claim.
\end{example}

% segment-id: o014.li2.chapter2.serre-subcategories-and-k0-groups.p006
We now introduce an important construction closely related to Serre
subcategories.

% segment-id: o014.li2.chapter2.serre-subcategories-and-k0-groups.d003
\begin{definition}\label{def:K0}
\index{K0 group@$\mathrm{K}_0$ group}
\index[sym1]{K0@$\mathrm{K}_0$}
Let $\mathcal{A}$ be an abelian category. Define the commutative group
$\mathrm{K}_0(\mathcal{A})$ as follows, writing the group operation
additively. Take the free $\Z$-module $F$ with basis $\Obj(\mathcal{A})$, and
write the element corresponding to $X\in\Obj(\mathcal{A})$ as
$\lrangle{X}\in F$. Let $R\subset F$ be the submodule generated by the
elements
% segment-id: o014.li2.chapter2.serre-subcategories-and-k0-groups.q002
\[\lrangle{X}-\lrangle{X'}-\lrangle{X''},\quad
0\to X'\to X\to X''\to0:\ \text{a short exact sequence in }\mathcal{A},\]
and call $\mathrm{K}_0(\mathcal{A}):=F/R$ the \emph{$\mathrm{K}_0$ group} of
$\mathcal{A}$. We henceforth denote the image of
$X\in\Obj(\mathcal{A})$ in $\mathrm{K}_0(\mathcal{A})$ by $[X]$.
\end{definition}

% segment-id: o014.li2.chapter2.serre-subcategories-and-k0-groups.lm001
\begin{lemma}\label{prop:K0-prep}
For every abelian category $\mathcal{A}$, the following equalities hold in
$\mathrm{K}_0(\mathcal{A})$.
% segment-id: o014.li2.chapter2.serre-subcategories-and-k0-groups.l003
\begin{compactenum}[(i)]
% segment-id: o014.li2.chapter2.serre-subcategories-and-k0-groups.i006
	\item $[0]=0$;
% segment-id: o014.li2.chapter2.serre-subcategories-and-k0-groups.i007
	\item if $X,Y\in\Obj(\mathcal{A})$ and $X\simeq Y$, then $[X]=[Y]$;
% segment-id: o014.li2.chapter2.serre-subcategories-and-k0-groups.i008
	\item $[X\oplus Y]=[X]+[Y]$;
% segment-id: o014.li2.chapter2.serre-subcategories-and-k0-groups.i009
	\item if
	$0\to X^1\xrightarrow{f^1}\cdots\xrightarrow{f^{n-1}}X^n\to0$
	is an exact sequence in $\mathcal{A}$, then
	$\sum_{i=1}^{n}(-1)^i[X^i]=0$;
% segment-id: o014.li2.chapter2.serre-subcategories-and-k0-groups.i010
	\item in the notation of \S\ref{sec:Abel-cat-subobjects}, for a family of
	subobjects $X=X_0\supset\cdots\supset X_r=0$, one has
% segment-id: o014.li2.chapter2.serre-subcategories-and-k0-groups.q003
	\[[X]=\sum_{i=0}^{r-1}[X_i/X_{i+1}].\]
\end{compactenum}
\end{lemma}

% segment-id: o014.li2.chapter2.serre-subcategories-and-k0-groups.pf004
\begin{proof}
Take the short exact sequence $0\to0\to0\to0\to0$; this gives (i). If
$X\simeq Y$, then $0\to X\rightiso Y\to0\to0$, together with (i), gives
(ii). The short exact sequence of Proposition~\ref{prop:biproduct-ses} gives
(iii).

For (iv), append zero terms and extend the exact sequence infinitely in both
directions. Consider the short exact sequences
% segment-id: o014.li2.chapter2.serre-subcategories-and-k0-groups.q004
\[0\to\Ker(f^i)\to X^i\to\Image(f^i)\to0,\quad i=1,\ldots,n.\]
Then
$\left[X^i\right]=\left[\Ker(f^i)\right]+\left[\Image(f^i)\right]
=\left[\Ker(f^i)\right]+\left[\Ker(f^{i+1})\right]$;
now take the alternating sum.

Finally, (v) follows by induction on $r$, using the short exact sequence
$0\to X_1\to X_0\to X_0/X_1\to0$.
\end{proof}

% segment-id: o014.li2.chapter2.serre-subcategories-and-k0-groups.rem001
\begin{remark}
\index{size issue}
Since this book requires groups to be realized on small sets, and in view of
Lemma~\ref{prop:K0-prep} (ii), a fully precise formulation of
$\mathrm{K}_0(\mathcal{A})$ should require $\mathcal{A}$ to have a small
skeleton \cite[Lemma 2.2.12]{Li1}. This issue has little effect here.
\end{remark}

% segment-id: o014.li2.chapter2.serre-subcategories-and-k0-groups.th002
\begin{theorem}[Euler--Poincaré principle]\label{prop:EP}
Let
$\cdots\to X^i\xrightarrow{d_X^i}X^{i+1}\to\cdots$ be a complex in an
abelian category $\mathcal{A}$ with only finitely many nonzero terms. Then the
equality
% segment-id: o014.li2.chapter2.serre-subcategories-and-k0-groups.q005
\[\sum_i(-1)^i\left[\Hm^i(X)\right]
=\sum_i(-1)^i\left[X^i\right],\quad
\Hm^i(X):=\Hm\left[X^{i-1}\xrightarrow{d_X^{i-1}}X^i
\xrightarrow{d_X^i}X^{i+1}\right]\]
holds in $\mathrm{K}_0(\mathcal{A})$.
\end{theorem}

% segment-id: o014.li2.chapter2.serre-subcategories-and-k0-groups.pf005
\begin{proof}
For every $i$ there is an exact sequence
% segment-id: o014.li2.chapter2.serre-subcategories-and-k0-groups.q006
\[0\to\Ker\left(d_X^{i-1}\right)\to X^{i-1}
\xrightarrow{d_X^{i-1}}\Ker\left(d_X^i\right)
\to\Hm^i(X)\to0,\]
and all terms are $0$ for $|i|\gg0$. Apply Lemma~\ref{prop:K0-prep} (iv) and
sum in $\mathrm{K}_0(\mathcal{A})$ to obtain
% segment-id: o014.li2.chapter2.serre-subcategories-and-k0-groups.q007
\[\sum_i(-1)^i
\left(\left[X^{i-1}\right]+\left[\Hm^i(X)\right]\right)
=\sum_i(-1)^i
\left(\left[\Ker\left(d_X^{i-1}\right)\right]
+\left[\Ker\left(d_X^i\right)\right]\right).\]
The terms on the right-hand side cancel in pairs; rearranging gives
$\sum_i(-1)^i[\Hm^i(X)]=\sum_i(-1)^i[X^i]$.
\end{proof}

% segment-id: o014.li2.chapter2.serre-subcategories-and-k0-groups.ex003
\begin{example}
Let $\mathcal{A}$ be the abelian category of finite-dimensional vector spaces
over a division ring $D$. Since every vector space has a basis,
$\mathrm{K}_0(\mathcal{A})=\Z\cdot[D]$. On the other hand, the assignment
$[X]\mapsto\dim_D X$ defines a group homomorphism
$\dim:\mathrm{K}_0(\mathcal{A})\to\Z$ sending $[D]$ to $1$. Thus
$\dim:\mathrm{K}_0(\mathcal{A})\rightiso\Z$.

If infinite-dimensional vector spaces are allowed, the $\mathrm{K}_0$ group
becomes trivial. Indeed, for every $D$-vector space $V$, a cardinality
argument gives a $D$-vector space $W$ such that $V\oplus W\simeq W$ and
$\dim_D W=\max\{\dim_D V,\aleph_0\}$; consequently, $[V]=0$.
\end{example}

% segment-id: o014.li2.chapter2.serre-subcategories-and-k0-groups.p007
For an exact functor $F:\mathcal{A}\to\mathcal{B}$ between abelian
categories, since $F$ preserves short exact sequences, the assignment
$[X]\mapsto[FX]$ defines a group homomorphism
\footnote{Editorial correction (O014-C028): the source prints
$\mathrm{K}_0(f)$, but this homomorphism is induced by the functor $F$; we use
$\mathrm{K}_0(F)$.}
$\mathrm{K}_0(F):\mathrm{K}_0(\mathcal{A})
\to\mathrm{K}_0(\mathcal{B})$.
Therefore, if $\mathcal{A}$ is an abelian subcategory of $\mathcal{B}$, there
is a natural homomorphism
$\mathrm{K}_0(\mathcal{A})\to\mathrm{K}_0(\mathcal{B})$.

% segment-id: o014.li2.chapter2.serre-subcategories-and-k0-groups.pr002
\begin{proposition}\label{prop:K0-exact-seq}
Let $\mathcal{T}$ be a Serre subcategory of an abelian category
$\mathcal{A}$, and write its Serre quotient as
$\mathcal{B}:=\mathcal{A}/\mathcal{T}$. Then the exact functors
$\mathcal{T}\to\mathcal{A}\xrightarrow{Q}\mathcal{B}$ induce an exact
sequence of additive groups
% segment-id: o014.li2.chapter2.serre-subcategories-and-k0-groups.q008
\[\mathrm{K}_0(\mathcal{T})\to\mathrm{K}_0(\mathcal{A})
\xrightarrow{\mathrm{K}_0(Q)}\mathrm{K}_0(\mathcal{B})\to0.\]
\end{proposition}

% segment-id: o014.li2.chapter2.serre-subcategories-and-k0-groups.pf006
\begin{proof}
First, $Q:\mathcal{A}\to\mathcal{B}$ is essentially surjective, so
$\mathrm{K}_0(\mathcal{A})\to\mathrm{K}_0(\mathcal{B})$ is surjective. The
composite $\mathrm{K}_0(\mathcal{T})\to\mathrm{K}_0(\mathcal{A})
\to\mathrm{K}_0(\mathcal{B})$ is plainly $0$. It remains to prove
$\Ker(\mathrm{K}_0(Q))\subset
\Image\left[\mathrm{K}_0(\mathcal{T})\to
\mathrm{K}_0(\mathcal{A})\right]$.

Let $\sum_{i=1}^n a_i[X_i]\in\Ker(\mathrm{K}_0(Q))$, where
$a_1,\ldots,a_n\in\Z$. This is equivalent to the existence of a family of
short exact sequences in $\mathcal{B}$,
$0\to Y'_j\xrightarrow{f_j}Y_j\xrightarrow{g_j}Y''_j\to0$, and integers
$b_j\in\Z$ for $j=1,\ldots,m$, such that
% segment-id: o014.li2.chapter2.serre-subcategories-and-k0-groups.q009
\begin{equation}\label{eqn:K0-exact-seq-aux-0}
\sum_{i=1}^n a_i\lrangle{QX_i}
=\sum_{j=1}^m b_j
\left(\lrangle{Y_j}-\lrangle{Y'_j}-\lrangle{Y''_j}\right)
\end{equation}
in the free $\Z$-module with basis $\Obj(\mathcal{B})$. Since
$Q:\mathcal{A}\to\mathcal{B}$ is constructed by localization in
Theorem~\ref{prop:Serre-quotient}, that theorem together with
\S\ref{sec:cat-localization} gives:
% segment-id: o014.li2.chapter2.serre-subcategories-and-k0-groups.l004
\begin{compactitem}
% segment-id: o014.li2.chapter2.serre-subcategories-and-k0-groups.i011
	\item $Q$ identifies $\Obj(\mathcal{A})$ with $\Obj(\mathcal{B})$, so
	\eqref{eqn:K0-exact-seq-aux-0} gives the equality in
	$\mathrm{K}_0(\mathcal{A})$
% segment-id: o014.li2.chapter2.serre-subcategories-and-k0-groups.q010
	\begin{equation}\label{eqn:K0-exact-seq-aux-1}
	\sum_{i=1}^n a_i[X_i]
	=\sum_{j=1}^m b_j\left([Y_j]-[Y'_j]-[Y''_j]\right);
	\end{equation}
% segment-id: o014.li2.chapter2.serre-subcategories-and-k0-groups.i012
	\item moreover, after replacing $Y_j$, $Y'_j$, $Y''_j$, and $f_j$ and $g_j$ by
	isomorphic objects and morphisms in $\mathcal{B}$ arising from $S$, we may
	arrange that there are morphisms $u_j$ and $v_j$ in $\mathcal{A}$ such that
	$f_j=Q(u_j)$ and $g_j=Q(v_j)$. From
	$Q(v_ju_j)=g_jf_j=0$ and Corollary~\ref{prop:fraction-zero}, a further
	adjustment by members of $S$ lets us arrange that $v_ju_j=0$.
\end{compactitem}

Thus, for every $1\leq j\leq m$, there are exact sequences in $\mathcal{A}$
% segment-id: o014.li2.chapter2.serre-subcategories-and-k0-groups.g005
\begin{gather*}
	0\to\Ker(v_j)\to Y_j\xrightarrow{v_j}Y''_j
	\to\Coker(v_j)\to0,\\
	0\to\Ker(u_j)\to Y'_j\xrightarrow{u_j}\Ker(v_j)
	\to\frac{\Ker(v_j)}{\Image(u_j)}\to0.
\end{gather*}
Since $Q$ is exact, $\Coker(v_j)$, $\Ker(u_j)$, and
$\frac{\Ker(v_j)}{\Image(u_j)}$ all belong to
$\mathcal{T}=\Ker(Q)$. Therefore, in $\mathrm{K}_0(\mathcal{A})$,
% segment-id: o014.li2.chapter2.serre-subcategories-and-k0-groups.q011
\[[Y_j]-[Y'_j]-[Y''_j]\in
\Image\left[\mathrm{K}_0(\mathcal{T})
\to\mathrm{K}_0(\mathcal{A})\right].\]
Substituting these relations into \eqref{eqn:K0-exact-seq-aux-1} gives
$\Ker(\mathrm{K}_0(Q))\subset
\Image\left[\mathrm{K}_0(\mathcal{T})
\to\mathrm{K}_0(\mathcal{A})\right]$.
\end{proof}

% segment-id: o014.li2.chapter2.serre-subcategories-and-k0-groups.p008
For an exact sequence such as the one in
Proposition~\ref{prop:K0-exact-seq}, a strategy that invariably succeeds is to
try to extend it to the left: seek a family of higher K-groups
$\mathrm{K}_i(\cdot)$ for $i\in\Z_{\geq0}$ together with a canonical exact
sequence
% segment-id: o014.li2.chapter2.serre-subcategories-and-k0-groups.q012
\[\cdots\to\mathrm{K}_{i+1}(\mathcal{B})
\to\mathrm{K}_i(\mathcal{T})\to\mathrm{K}_i(\mathcal{A})
\to\mathrm{K}_i(\mathcal{B})\to\cdots\]
One also hopes that the family
$(\mathrm{K}_i(\mathcal{A}))_{i\geq0}$ contains deep information about
$\mathcal{A}$ and is, at least to some extent, computable. This is the subject
of K-theory, whose applications also encompass exact categories, which are
more general than abelian categories. Because the relevant constructions are
based on the viewpoint of homotopy theory, this book cannot treat them in
detail; interested readers may consult \cite{Lai19}, or the summary in
\cite[\S 13.6]{Bu10}.
\index{exact category}
